% The security and privacy review, docs/security.md in full: executive summary, scope, trust
% boundaries, assets, threat model, authentication, authorization, Privacy Mode and the gate
% (controls, file rules, fail closed, pseudonyms, verification), data at rest and in transit,
% web application security, untrusted input, audit logging, background jobs and plug-ins, supply
% chain, regulatory mapping, the deployment checklist, limitations, verification.
% The code is the authority (medorama/src/server.rs, signin.rs, users.rs, gate.rs, deid.rs,
% dicom.rs, pseudonyms.rs, gateway.rs, audit.rs, actions.rs, db.rs, sql/pg.rs, config.rs).
\chapter{Security and privacy review}
\label{chap:security}

This chapter is for the people who decide whether Medorama may run on their institution's
patient data: IT security reviewers, privacy officers, and IT and research administrators. It
describes what Medorama protects, against whom, how, and what it leaves to the deploying
institution. Each control names the code that implements it and the test that checks it, so that
every claim can be verified rather than taken on trust. It follows the repository's
\file{docs/security.md}.

The chapter describes Medorama \mdversion{} (unreleased). Medorama is research software, not a
medical device, and it holds no certification. To report a vulnerability, follow the
repository's \file{SECURITY.md}. The system's parts are described in \cref{chap:architecture}, a
plain-language explanation of how identifiers are hidden is \cref{chap:protecting}, and sign-in
is described in \cref{chap:sign-in}.

\section{Executive summary}
\label{sec:security-summary}

Medorama is a single server program, written in Rust, that an institution runs on its own hosts.
It serves a web page and three viewers to the institution's users, and reads patient data from
imported tables, files and connected databases. It sends nothing to any third party: no
telemetry, no external services, and no scripts or fonts from other sites.
\begin{itemize}
\item \textbf{One way in.} Every request passes Medorama's sign-in. Users sign in either with
  Medorama's own accounts (Argon2id password hashes, throttling, hashed session tokens) or
  through the institution's authenticating reverse proxy, which brings single sign-on and
  multi-factor authentication. The sign-in method is one module (\file{signin.rs}); roles and
  permissions are decided in Medorama, on the server, for every request.
\item \textbf{Identifiers in one place.} Names, medical record numbers, dates of birth, account
  numbers and imported IDs are kept only in the catalog database, in their own tables. The
  viewers, indexes, results and plug-ins receive pseudonyms only.
\item \textbf{Privacy Mode by default.} Users without the PHI permission are always in Privacy
  Mode, and users with it start in Privacy Mode unless they are administrators. In Privacy Mode,
  the privacy gate de-identifies every answer as it is sent: identifiers become placeholders and
  random pseudonyms, dates move back by a per-patient offset, and files are de-identified at the
  source (WFDB, GE MUSE XML, DICOM with the actions of PS3.15's Basic Profile for the
  attributes it lists, EDF and BDF). Pictures,
  videos and PDF documents pass with a warning; any other file is refused.
\item \textbf{Fail closed.} A last layer of the server puts every JSON answer of a Privacy Mode
  session through the gate and refuses any other answer that the gate did not clear, so a new
  endpoint that forgets the gate cannot leak identifiers. An end-to-end test seeds identifiers in
  every format and checks that no byte of any answer, file or export contains them.
\item \textbf{Read-only sources, sealed audit log.} Data folders and connected databases are only
  read. Sign-ins, identifier searches and views, exports and administrative changes are written to
  an audit log that is sealed every minute with a hash chain and can be copied to storage out of
  the database administrators' reach.
\item \textbf{Conservative network posture.} Medorama listens on 127.0.0.1 by default and relies
  on a TLS reverse proxy or an SSH tunnel for remote access. Its waveform viewer runs on a private
  loopback port behind an allow-list gateway. Connections to PostgreSQL use verified TLS.
\end{itemize}
The deploying institution remains responsible for the host, TLS, backups, accounts and the PHI
permission, network policy, and the governance of exports. The main residual risks are identities
burned into images and free text, quasi-identifiers, the trust placed in administrators and
installed plug-ins, and, in proxy mode, the trust placed in the proxy
(\cref{sec:security-limitations}).

\section{Scope and intended use}
\label{sec:security-scope}

Medorama brings a patient's data from many databases together: admissions and bed stays,
laboratory results, vital signs, medications, conditions, notes, ECGs, imaging and bedside-monitor
recordings. It shows each patient's journey on one page, lets researchers build cohorts across
every database, and opens recordings and images in its own viewers.
\begin{itemize}
\item \textbf{Research use.} Medorama is research software. It is not a medical device, is not
  cleared for diagnosis or treatment, and must not be used to make clinical decisions.
\item \textbf{One institution.} It runs on hosts that the institution controls: one server with a
  SQLite catalog, or several servers sharing a PostgreSQL catalog. It serves that institution's
  users through a browser. It is not a multi-tenant cloud service.
\item \textbf{Data enter as databases.} An administrator imports tables (CSV, Excel, Parquet)
  exported from the institution's systems. Files (recordings, ECGs, DICOM studies, documents)
  enter only as rows of such tables, which name their patient; they are read where they lie and
  never written (\cref{chap:importing}).
\item \textbf{Or they are read in place.} An administrator may connect databases (an OMOP CDM
  instance, a warehouse), which Medorama reads read-only at each request through views of fixed
  columns (\cref{chap:connecting}). Only the identifiers of their persons are copied, when the
  persons are linked to patients.
\end{itemize}
What Medorama does \textbf{not} provide, and a deployment must: TLS termination, the identity
provider for single sign-on or multi-factor authentication, host hardening, backups and their
encryption, network policy, and a retention policy (\cref{sec:security-checklist}).

\section{Architecture and trust boundaries}
\label{sec:security-boundaries}

The system context and its trust boundaries are drawn in \cref{fig:architecture-context}, and
the path of a request through the server in \cref{fig:architecture-components}.
\Cref{tab:security-zones} lists the zones and how far each is trusted.

\begin{xltabular}{\linewidth}{@{}P{2.8cm}P{5.0cm}L@{}}
\caption{The zones and how far each is trusted.}
\label{tab:security-zones}\\
\toprule
\tabhead Zone & \tabhead Contains & \tabhead Trust \\
\midrule
\endfirsthead
\toprule
\tabhead Zone & \tabhead Contains & \tabhead Trust \\
\midrule
\endhead
\bottomrule
\endfoot
User's computer & The browser, with the Medorama page and the ECG, PACS and VitalStream viewers
  (static pages) & \textbf{Untrusted.} Every request is authenticated and authorized on the
  server; the page decides nothing about access. \\
Network & HTTPS from a reverse proxy, or an SSH tunnel & \textbf{Untrusted.} Medorama listens on
  127.0.0.1 by default. \\
Authenticating proxy (optional) & The institution's sign-in in front of Medorama
  (\code{sign\_in = "proxy"}) & \textbf{Trusted to name the user}, on connections from
  \code{proxy\_addresses} only (\cref{sec:security-proxy}). \\
Server host & The \cmd{medorama} process; its VitalStream viewer (a child process on a private
  loopback port); plug-in actions (child processes); the catalog database; the index and results
  folders & \textbf{Trusted.} One operating-system account, the catalog's owner, runs
  everything. \\
Data sources & Exported tables, recordings, ECGs, DICOM studies, documents & Read only, inside
  the configured \code{data\_roots}. \\
Database server (optional) & The catalog in PostgreSQL, when \code{database\_url} is set &
  \textbf{Trusted like the host}: it holds every identifier. It is reached over TLS with a
  verified certificate, by an account with rights on that database only. \\
Connected databases (optional) & OMOP instances, warehouses & \textbf{Read only, enforced}
  (read-only sessions, read-only SQLite files). What they return is data, never instructions:
  only the canonical columns are read. \\
\end{xltabular}

Every HTTP request passes the same layers before and after its handler (\file{server.rs},
\code{router}):
\begin{enumerate}
\item \textbf{the guard} (\code{guard}) sets security headers on every answer, and refuses any
  API request that changes something unless it carries the header \code{X-Medorama: 1};
\item \textbf{the privacy layer} (\code{privacy}) identifies the user once per request through
  \code{signin::identify}, decides whether the session is in Privacy Mode, and on the way out
  applies the privacy gate to every answer of such a session (\cref{sec:security-privacy});
\item \textbf{the handler} receives the user as \code{Auth(User)} and checks the role and the PHI
  permission it needs.
\end{enumerate}
The deployment topologies (one server, several servers on PostgreSQL, behind a TLS proxy or an
institution's sign-in, viewers alone) are described in \cref{sec:architecture-topologies}. Every
setting and endpoint is specified in \cref{chap:settings-reference} and \cref{chap:http-api}.

\section{Assets and data classification}
\label{sec:security-assets}

\begin{xltabular}{\linewidth}{@{}P{3.0cm}P{3.0cm}L@{}}
\caption{What there is to protect, how it is classified, and where it lives.}
\label{tab:security-assets}\\
\toprule
\tabhead Asset & \tabhead Classification & \tabhead Where it lives \\
\midrule
\endfirsthead
\toprule
\tabhead Asset & \tabhead Classification & \tabhead Where it lives \\
\midrule
\endhead
\bottomrule
\endfoot
Direct identifiers: names, MRNs, dates of birth, account numbers (CSN), imported patient and
  encounter IDs & \textbf{PHI} & The catalog's identifier tables (\code{patient\_identifiers},
  \code{patient\_mrns}, \code{encounter\_identifiers}) and the rows of imported databases
  (\code{import\_rows}) \\
Clinical data: results, vital signs, medications, conditions, notes, with their real dates &
  \textbf{PHI} (dates and free text identify; HIPAA counts them) & The catalog
  (\code{ehr\_items}, \code{import\_rows}) \\
Files: recordings, ECGs, DICOM studies, PDFs, pictures, videos & \textbf{PHI} (headers, file
  names, burned-in text) & The data sources (read only) \\
Derived data: recording indexes, derived series, plug-ins' results & Sensitive (dated,
  linkable) & The index and results folders \\
Pseudonyms (\code{MO-XXXX-XXXX}, \code{MO-E-XXXX-XXXX}) & Pseudonymous; the link to identity
  exists only in the catalog & The catalog \\
Credentials: password hashes, session token hashes & Secret & The catalog (\code{users},
  \code{sessions}) \\
Database passwords: the PostgreSQL catalog's, each connected database's & Secret & Files
  readable by the catalog's account (\code{database\_password\_file}, a connection's
  \code{password\_file}) or the environment; never in the database, a settings file, a URL or
  the audit log \\
Connected persons: each connected database's own key for a linked person & \textbf{PHI} (a
  source's key may identify) & \code{source\_persons}; never sent: persons leave as pseudonyms,
  entries as tokens \\
Secrets: the catalog's secret (it keys the replacement of DICOM UIDs and the tokens of connected
  entries); the internal viewer's password file and signing key; the database key (encrypted
  builds) & Secret & The database (\code{settings}), \file{<home>/viewer/auth.json} (mode 0600),
  the environment or a key file \\
The audit log: who signed in, imported, searched identifiers, viewed, exported & Sensitive
  (searches may name patients) & The catalog (\code{audit}, \code{audit\_seals}), and optionally
  \code{audit\_export} \\
\end{xltabular}

\section{Threat model}
\label{sec:security-threats}

\Cref{tab:security-actors} lists the actors considered and the main controls against each, and
\cref{tab:security-stride} the threats by STRIDE category, their mitigations and what remains.

\begin{xltabular}{\linewidth}{@{}P{3.4cm}P{3.6cm}L@{}}
\caption{The actors considered, and the main controls against each.}
\label{tab:security-actors}\\
\toprule
\tabhead Actor & \tabhead Capability & \tabhead Main controls \\
\midrule
\endfirsthead
\toprule
\tabhead Actor & \tabhead Capability & \tabhead Main controls \\
\midrule
\endhead
\bottomrule
\endfoot
Network attacker & Reaches the server's port & Loopback by default; TLS at a reverse proxy;
  every endpoint except sign-in and the static pages requires a signed-in user; sign-in
  throttling \\
Someone who can reach Medorama's port in proxy mode & Sends a forged user header & The header is
  believed only from \code{proxy\_addresses}; Medorama stays on 127.0.0.1 or behind a firewall
  (\cref{sec:security-proxy}) \\
Signed-in user without the PHI permission & Any role & Always in Privacy Mode: answers pass
  through the gate; identifier searches refused; files the gate cannot de-identify refused \\
User with the PHI permission & Sees identifiers when they choose & Privacy Mode on by default
  for non-administrators; switching it off is audited; identifier searches, views and exports are
  audited \\
Administrator & Manages users, imports and settings & Trusted by design; actions audited; the
  audit log sealed and exportable \\
Malicious or malformed file & In a data source & Parsers check every length; failures are
  errors, not crashes; active content sandboxed \\
Malicious web content & A crafted SVG, PDF or note & A strict Content Security Policy, no inline
  scripts, sandboxed file responses, output escaping \\
Other users of the host & Local accounts & Owner-only database and secret files; the internal
  viewer requires a password only Medorama knows. In proxy mode, local accounts must be trusted
  (\cref{sec:security-limitations}) \\
Lost or stolen storage & A disk, a backup & Optional database encryption (SQLCipher); full-disk
  encryption and encrypted backups are the deployment's responsibility \\
A connected database returning crafted data & Values meant to break parsing, the page or the
  gate & Only the canonical columns are read, as values; at most \code{max\_rows} rows per
  criterion, each query time-limited; everything shown is escaped and passes the gate \\
Network between Medorama and its databases & Reads or alters traffic & TLS (rustls) with the
  server's certificate and name verified; required for any server not on the same host unless
  the URL says \code{sslmode=disable} \\
\end{xltabular}

\begin{xltabular}{\linewidth}{@{}P{2.3cm}P{3.2cm}P{4.6cm}L@{}}
\caption{Threats (STRIDE), their mitigations, and the residual risk.}
\label{tab:security-stride}\\
\toprule
\tabhead Threat & \tabhead Example & \tabhead Mitigation & \tabhead Residual risk \\
\midrule
\endfirsthead
\toprule
\tabhead Threat & \tabhead Example & \tabhead Mitigation & \tabhead Residual risk \\
\midrule
\endhead
\bottomrule
\endfoot
Spoofing & Guessing a password, stealing a session, forging the proxy's header & Argon2id
  hashes; 256-bit random session tokens stored hashed; \code{HttpOnly}, \code{SameSite=Strict}
  cookies, \code{Secure} behind TLS; throttling. In proxy mode, single sign-on and MFA at the
  proxy, and the header believed only from \code{proxy\_addresses} & Password mode has no MFA
  (use proxy mode where it is required). Proxy mode is as strong as the proxy and the network
  path to Medorama \\
Tampering & Changing data through forged requests & Changing requests need \code{X-Medorama: 1}
  (CSRF) and, in password mode, a \code{SameSite=Strict} cookie; roles checked on the server &
  An administrator can change anything, by design; audited \\
Repudiation & Denying a search or an export & An audit log of sign-ins, imports, removals,
  identifier searches and views, exports and Privacy Mode switches, sealed every minute with a
  hash chain and exportable to a file out of the database's reach & Changing the database and
  the export file together \\
Information disclosure & A user without the PHI permission reading identifiers & The privacy
  gate on every answer, failing closed; identifier searches refused; pseudonyms not derivable
  from identifiers & Pixels, PDFs, videos and free text may still show identity: they are marked,
  not removed \\
Denial of service & Huge uploads, expensive searches & Uploads capped at 1\,GiB; imports for
  administrators only; a worker pool; time and row limits on connected databases & Not hardened
  against hostile signed-in users \\
Elevation of privilege & A viewer reaching outside its files; path traversal & File access only
  through an item's own files inside \code{data\_roots}; the gateway's allow-list; \code{..}
  refused & Plug-ins run with the catalog's rights (administrator-installed code is trusted) \\
\end{xltabular}

\section{Authentication}
\label{sec:security-auth}

All sign-in code is in \file{signin.rs}, and accounts and sessions are in \file{users.rs}. The
setting \code{sign\_in} chooses the method. The configuration, proxy examples and the procedure
are in \cref{chap:sign-in}.

\subsection{Password sign-in (the default)}
\label{sec:security-password}

\begin{itemize}
\item \textbf{Accounts} are personal, created by an administrator in \ui{Admin} $\rightarrow$
  \ui{Users} or with \cmd{medorama user add}. Account names use letters, digits and
  \code{. \_ - @}, up to 64 characters.
\item \textbf{Passwords} have at least 10 characters and are stored as \textbf{Argon2id} hashes
  (the \code{argon2} crate's defaults: 19\,MiB of memory, 2 iterations, 1 lane, a random salt
  per password). A sign-in with an unknown name performs the same hashing work, so the timing
  does not reveal which names exist.
\item \textbf{Throttling.} Every attempt is counted before its password is checked, so attempts
  made at once count too. After 5 attempts that failed within 60 seconds, an account name is
  locked for 60 seconds \textbf{from that address} (the lock is per name and address, so nobody
  can lock another user out from elsewhere); after 30 failed attempts within 60 seconds from one
  address, whatever the names, that address is refused for 60 seconds. Behind a reverse proxy in
  password mode every request comes from the proxy's address, so these limits are the whole
  site's. At most 4 passwords are checked at once (Argon2id takes memory on purpose); a name no
  account can have (over 64 characters, other characters) is refused without a check. Every
  failure waits one second before the answer. Failures are audited (\code{login\_failed}, with
  at most 64 characters of the name given).
\item \textbf{Sessions.} A session is a 256-bit random token from the operating system's
  generator, and only its SHA-256 is stored. The cookie (\code{medorama\_<port>}) is
  \code{HttpOnly}, \code{SameSite=Strict} and \code{Path=/}, and \code{Secure} when
  \code{secure\_cookies = true}, which is required behind TLS. Sessions expire after
  \code{session\_hours} (default 12) and end when unused for \code{session\_idle\_minutes}
  (default 60; with 0, only \code{session\_hours} ends them). The page's automatic checks for new
  data (\code{GET version}, and requests marked \code{X-Background: 1}) do not count as use, so
  an unattended page signs out. The time of last use is written at most once a minute, and
  expired sessions are purged.
\item \textbf{Ending sessions.} \ui{Sign out} deletes the session. A user can end all their
  sessions on every device (\ui{Sign out everywhere}, \code{POST logout \{"everywhere":
  true\}}), and an administrator can end a user's (\ui{Admin} $\rightarrow$ \ui{Users}
  $\rightarrow$ \ui{Sign out}, \code{PATCH users/\{id\} \{"sign\_out": true\}}); both are
  audited. A new password or a disabled account ends them too, and a disabled account is refused
  at its next request.
\end{itemize}
Tests: \file{medorama/tests/signin.rs} (\code{with\_passwords}), and the sign-in steps of every
API test.

\subsection{Sign-in through the institution's proxy}
\label{sec:security-proxy}

With \code{sign\_in = "proxy"}, an authenticating reverse proxy (oauth2-proxy, Apache httpd with
mod\_auth\_openidc, mod\_shib or mod\_auth\_gssapi, nginx with \code{auth\_request}) signs the
user in with the institution's method, including single sign-on (OpenID Connect, SAML), Kerberos
and multi-factor authentication, and names the user in a header (\code{proxy\_user\_header},
default \code{X-Remote-User}).
\begin{itemize}
\item \textbf{The header is believed only on connections from \code{proxy\_addresses}} (default
  127.0.0.1 and ::1; exact addresses, no ranges). From anywhere else it is ignored, and the
  request is not signed in (\code{signin::identify}).
\item \textbf{Accounts stay in Medorama.} The name is matched without regard to case with an
  enabled account (\code{users::get\_by\_login}). A name with no account, a disabled account, or
  two accounts that differ only in case is refused with HTTP 403 and a reason that the page
  displays. Roles and the PHI permission are assigned in Medorama, not taken from the identity
  provider. An account created for proxy sign-in needs no password: it receives a random one that
  nobody knows (\cref{sec:sign-in-accounts}).
\item \textbf{No password is accepted.} \code{POST login} and password changes are refused (HTTP
  400).
\item \textbf{Sessions are the proxy's.} Medorama keeps none in this mode, and \ui{Sign out}
  leads to \code{proxy\_logout\_url}. Session lifetime and idle timeout are set at the proxy or
  at the identity provider.
\item \textbf{Audit.} One \code{login} entry is written per visit (detail \code{proxy}), and one
  \code{login\_failed} entry per visit for an unknown name.
\item \textbf{Requirements.} Medorama must be reachable only through the proxy (keep
  \code{host = "127.0.0.1"}, or firewall its port to the proxy's address), its host must have no
  untrusted local accounts, the proxy must remove any incoming copy of the header, and TLS must
  terminate at the proxy, with \code{secure\_cookies = true} (\cref{sec:sign-in-requirements}).
\item \textbf{Validation.} Inconsistent settings stop Medorama at start (\file{config.rs},
  \code{sign\_in\_method}).
\end{itemize}
The proxy's header must come once: a request with two copies of it (one the client sent, one the
proxy added) names nobody. Addresses compare as the same address whether written as IPv4 or as
IPv6's form of it (\code{::ffff:127.0.0.1}).

Tests: \file{medorama/tests/signin.rs} (\code{through\_a\_proxy}: the header from another address
counts for nothing, a name without an account is refused with its reason, passwords are refused,
one audit entry is written per visit; \code{proxy\_headers\_and\_addresses};
\code{throttled\_sign\_in}: 40 attempts at once, 5 checked;
\code{the\_phi\_permission\_is\_given\_by\_those\_who\_have\_it}) and the unit test
\code{sign\_in\_settings} of \file{config.rs}.

\subsection{Other access paths}
\label{sec:security-other-access}

\begin{itemize}
\item \textbf{The command line} (\cmd{medorama \dots} on the host) acts as the catalog's owner.
  Whoever can run it under the catalog's operating-system account can read the database directly
  anyway.
\item \textbf{The internal VitalStream viewer} listens on 127.0.0.1 on a private port, with a
  password file (\file{<home>/viewer/auth.json}, mode 0600) whose signing key only Medorama knows.
  Medorama's gateway signs each request it passes on and never forwards the user's own cookie.
  Users never sign in to the viewer, and its port is of no use to other local accounts
  (\file{viewer.rs}, \file{gateway.rs}).
\item \textbf{The viewers alone}, without Medorama, have their own optional password and a root
  folder fixed when they are built (\cref{sec:sign-in-viewers}).
\end{itemize}
Medorama itself does not provide multi-factor authentication or single sign-on (they are provided
through the proxy), password expiry, or group and role mapping from an identity provider.

\section{Authorization}
\label{sec:security-authorization}

\begin{xltabular}{\linewidth}{@{}P{2.2cm}L@{}}
\caption{The roles.}
\label{tab:security-roles}\\
\toprule
\tabhead Role & \tabhead May \\
\midrule
\endfirsthead
\toprule
\tabhead Role & \tabhead May \\
\midrule
\endhead
\bottomrule
\endfoot
\code{viewer} & Look at patients, encounters, recordings and databases; build and keep cohorts;
  keep a history. \\
\code{analyst} & Everything a viewer may do, and run actions (indexing, plug-ins). \\
\code{admin} & Everything, including importing and removing databases, managing users and
  settings, reviewing identities and reading the audit log. \\
\end{xltabular}

Separately from the role, the \textbf{PHI permission} decides whether a user may leave Privacy
Mode, that is, see and search identifiers. An administrator sets it when creating the account or
later.
\begin{itemize}
\item \textbf{Every check is on the server.} The page hides what a user may not do, but each
  endpoint checks again (\code{admin()}, \code{User::can\_run}, the session's mode).
\item \textbf{A session in Privacy Mode has no PHI permission.} The server replaces its user with
  a copy without the permission (\code{User::in\_privacy\_mode}), so any check that depends on
  the permission fails closed for that session.
\item \textbf{Cohorts} are private to their owner unless shared, to read or to edit, with chosen
  users; only the owner shares or deletes a cohort (\cref{chap:cohorts}).
\item \textbf{Files} are served one item at a time (\code{patients/\{id\}/file/\{item\}}) and, by
  name, only the files that belong to that item (a WFDB record's signal files, an attached
  folder's files), inside \code{data\_roots}; never other files of the folders they lie in. Paths
  are compared with \code{data\_roots} by their real paths (symbolic links resolved); an attached
  folder's symbolic links are not its files, and a file named by a WFDB header must really lie
  beside it.
\item \textbf{The PHI permission} is given (to a new account or an existing one, the
  administrator's own included) only by an administrator who has it.
\item \textbf{The VitalStream gateway} passes only the listed requests of the viewer's API, each
  checked against the user's rights: default deny (\cref{sec:architecture-gateway}). It is tested
  in \file{medorama/tests/gateway.rs}.
\item \textbf{Not provided:} per-patient or per-study access lists. Every signed-in user may see
  every patient in the catalog, as Privacy Mode and their permission allow. Separate catalogs
  serve separate studies.
\end{itemize}

\section{Privacy Mode and the privacy gate}
\label{sec:security-privacy}

The design is explained in full, for every audience, in \cref{chap:protecting}, and drawn in
\cref{fig:protecting-gate}. This section summarizes it as controls for a review.

\subsection{The rule}
\label{sec:security-rule}

\textbf{The catalog keeps every identifier as imported, and the server decides, at the moment of
sending, what leaves. In Privacy Mode, no identifier reaches the browser.} The page never
receives what it must not show and does no masking of its own. Users without the PHI permission
are always in Privacy Mode; users with it start in Privacy Mode and may leave it, except
administrators, who start outside it and may switch it on (\cref{tab:protecting-who}). The switch
calls \code{POST privacy}, which is audited (\code{privacy\_on}, \code{privacy\_off}). The choice
is the session's and is kept by the server, never by the browser: in password mode in the
session's row, until the session ends; in proxy mode for the user's visit (until an idle time
passes without a request, or \code{session\_hours}), in the memory of the server that was asked
(with several servers, another one answers in the user's default mode). Every sign-in starts in
the user's default mode, so no cookie, left by an earlier user of the browser or made by hand, can
take a session out of Privacy Mode. A user without the permission who asks to leave Privacy Mode
is refused (HTTP 403).

\subsection{Controls}
\label{sec:security-controls}

\begin{xltabular}{\linewidth}{@{}P{2.4cm}P{4.6cm}P{3.6cm}L@{}}
\caption{The controls of Privacy Mode, with their code and their tests.}
\label{tab:security-controls}\\
\toprule
\tabhead Control & \tabhead What it does & \tabhead Code & \tabhead Test \\
\midrule
\endfirsthead
\toprule
\tabhead Control & \tabhead What it does & \tabhead Code & \tabhead Test \\
\midrule
\endhead
\bottomrule
\endfoot
Mode decision & Privacy Mode unless the user has the PHI permission and switched it off &
  \file{server.rs} \code{in\_privacy}, \code{session\_user} & \file{tests/privacy.rs} \\
Identifier dictionary & Imported patient and encounter IDs become pseudonyms; MRNs and CSNs
  placeholders; recordings' paths tokens and the folders they lie in \code{[Folder]}, in every
  string of an answer (paths inside text too: an error that names a file). Rebuilt whenever an
  identifier changes (triggers on every table it is built from move
  \code{data\_version.identifiers}) & \file{gate.rs} \code{Dict}, \code{Restricted::json};
  \file{server.rs} \code{restricted} & \file{gate.rs} unit tests; \file{tests/privacy.rs}
  (\code{the\_dictionary\_follows\_the\_identifiers}) \\
Placeholders & Names, dates of birth, staff names, withheld and unrecognized columns become fixed
  text & \file{gate.rs} \code{column\_placeholder}, \code{column\_mask}, \code{cell} &
  \file{tests/privacy.rs} \\
Pseudonyms & Random, stored, with a check character; never derived from an identifier &
  \file{pseudonyms.rs} & \file{pseudonyms.rs} unit tests \\
Dates and ages & Every date moved back by the catalog's offset and the patient's own days; ages
  above 89 as \code{90+} & \file{deid.rs} \code{Settings}, \code{answer\_with},
  \code{shift\_text} & \file{deid.rs} unit tests; \file{tests/privacy.rs} (no JSON number is a
  real time) \\
Notes & The patient's own names (each part of a compound name: \code{Smith-Jones},
  \code{O'Brien}) and date of birth replaced; marked; withheld if the administrator chooses &
  \file{gate.rs} \code{Restricted::note}, \code{scrub\_names}, \code{name\_words} &
  \file{tests/privacy.rs} \\
Errors and jobs & The viewer's error texts replaced by a fixed sentence; the records of a folder
  it left out named by number; a failed job's error \code{failed}; results folders for
  administrators only (and not in Privacy Mode) & \file{gateway.rs} \code{private\_answer};
  \file{server.rs} \code{jobs\_list} & \file{tests/privacy.rs} (\code{around\_the\_gate},
  \code{the\_viewers\_errors\_in\_privacy\_mode}) \\
EDF+ annotations in VitalStream & Treated like notes: \code{[Annotation]} when notes are
  withheld, otherwise the patient's names and date of birth replaced & \file{gateway.rs}
  \code{gate\_annotations} & \file{gateway.rs} unit test \\
Searches & Searches of identifiers refused; patients found by pseudonym; every search, filter
  and sort uses only what the session is shown & \file{search.rs}, \file{query.rs},
  \file{patient/list.rs} & \file{tests/privacy.rs}; \file{tests/privacy\_search.rs} \\
Files & De-identified by kind, passed with a warning, or refused (\cref{sec:security-files}) &
  \file{gate.rs} \code{file\_kind}; \file{server.rs} \code{attachment\_response} &
  \file{tests/privacy.rs}; \file{tests/dicom\_samples.rs} \\
The VitalStream gateway & Recordings by token, no folder browsing, times translated both ways &
  \file{gateway.rs} & \file{tests/gateway.rs} \\
Every route & Every GET route of the API, and the POST routes that answer with data, asked in
  Privacy Mode; a new GET route fails the test until it has a case & & \file{tests/privacy.rs}
  (\code{every\_route\_in\_privacy\_mode}) \\
Exports and plug-ins & Exports pass the same gate; plug-ins get pseudonyms only &
  \file{cohorts/export.rs}, \file{databases.rs}; \file{jobs.rs} & \file{tests/privacy.rs} \\
Connected databases & Only canonical columns, each with a privacy class; persons as pseudonyms;
  entries as keyed tokens & \file{sources/}, \file{patient/items.rs} \code{live\_items} &
  \file{tests/connections.rs} \\
Fail closed & The last layer gates every JSON answer and refuses any other uncleared answer &
  \file{server.rs} \code{privacy} & \file{tests/privacy.rs} \\
\end{xltabular}

Administrators choose what Privacy Mode withholds, the catalog's date offset and the spread of
per-patient days in \ui{Admin} $\rightarrow$ \ui{Privacy Mode}, and each change is audited
(\code{deid\_settings}). The defaults follow the Safe Harbor identifiers that the catalog holds,
except the text of notes, which is shown with protections and a warning unless the administrator
withholds it. Pseudonyms, file tokens, staff placeholders, the withholding of unrecognized columns
and the refusal of unknown files always apply.

\subsection{The gate's file rules}
\label{sec:security-files}

\begin{xltabular}{\linewidth}{@{}P{3.0cm}L@{}}
\caption{The gate's rules for files in Privacy Mode.}
\label{tab:security-files}\\
\toprule
\tabhead Kind & \tabhead In Privacy Mode \\
\midrule
\endfirsthead
\toprule
\tabhead Kind & \tabhead In Privacy Mode \\
\midrule
\endhead
\bottomrule
\endfoot
WFDB header & Identifying comments removed (their keys compared without case, spaces, \code{\_}
  or \code{-}, and written \code{key: value}, \code{key=value} or \code{key - value}); the
  patient's own names and date of birth replaced in every other comment; staff
  \code{[Staff Name]}; dates moved; ages capped; file names replaced by tokens \\
WFDB samples & Sent as they are (numbers only), when a header of the same attachment names the
  file; any other \file{.dat} file refused \\
WFDB samples in a MATLAB file & When a header of the same attachment names the \file{.mat} file:
  its first matrix only (other variables after it not sent), its descriptive text and the
  matrix's name blanked (\code{mat\_deidentified} in
  \file{libs/engine/src/sources/wfdb\_files.rs}); a file that is not a level 4 or uncompressed
  level 5 MAT file, and any other \file{.mat} file, refused \\
GE MUSE XML & Demographics removed except age and sex (their tags found with attributes and
  spaces too); the order removed; staff replaced; the patient's own names and date of birth
  replaced in the statements; dates moved. A file that still holds a patient's name, ID or date
  of birth afterwards, and other XML, refused \\
DICOM & Header rewritten with the actions of PS3.15 Table E.1-1 (the Basic Application Level
  Confidentiality Profile) for the attributes \file{dicom.rs} lists, with Retain Longitudinal
  Temporal Information with Modified Dates; the free text left (and, in implicit-VR files, text
  of unknown VR) without the patient's names, date of birth and the catalog's identifiers; pixel
  data copied byte for byte; marked as possibly identifying. Not a certified implementation of
  the profile (\cref{sec:protecting-dicom}). Deflated and big-endian files, and sequences nested
  more than 32 deep, refused \\
EDF, EDF+, BDF, BDF+ & Patient field \code{X X X X}; recording field \code{Startdate <moved
  date> X X X} (EDF+, BDF+) or blank (EDF, BDF); start date and time moved back; annotation texts
  removed, time-keeping kept; samples as they are. EDF+ files with annotations larger than
  64\,MB refused \\
Pictures, videos, PDF & Sent as they are, marked \code{x-medorama-may-contain-phi} (an amber
  warning triangle on the viewer's tab) \\
Anything else & Refused (HTTP 403) \\
\end{xltabular}

In Privacy Mode, files are named by tokens (\code{mo-f0000.hea}) and real file names are not
accepted, and file answers carry \code{Vary: Cookie}. The element-by-element DICOM rules and the
EDF rules are in \cref{sec:protecting-dicom} and \cref{sec:protecting-edf}.

\subsection{Fail closed}
\label{sec:security-fail-closed}

The gate is applied by the handlers that know their data, and \textbf{enforced by the last layer}
(\file{server.rs}, \code{privacy}). For a session in Privacy Mode, every JSON answer passes the
gate there; a file passes only if a handler cleared it (with the internal header
\code{x-medorama-gate}, which is removed before the answer leaves); and any other answer carrying
data is refused (HTTP 403) and logged on the server's standard error. The pages, their scripts,
redirects, errors without data and the session's own endpoints (\code{me}, \code{login},
\code{logout}, \code{privacy}, \code{sign-in}) pass unchanged. A new endpoint that forgets the
gate therefore cannot leak: its answer is de-identified or refused.

\subsection{Pseudonyms and date shifting}
\label{sec:security-pseudonyms}

Pseudonyms are drawn from the operating system's secure random generator and stored. They are not
a hash of any identifier, which, for short MRNs, could be reversed by enumeration. This follows
HIPAA \S164.514(c) and ISO/TS 25237. Dates are shifted, not removed, because intervals are the
substance of research. The per-patient part of the offset (0 to \code{spread\_days}, default 365,
at most 3650) is drawn once and stored, so every session sees the same dates, and a patient's
intervals are exact. The offset never reaches the page of a user who is not an administrator
(\code{GET deid} is for administrators only, and \code{GET config} leaves it out), and dates
typed in Privacy Mode are compared, on the server, with each patient's dates as shown. Date
shifting is pseudonymization: learning one real date of a patient reveals that patient's
offset, and every other patient's to within \code{spread\_days}.

\subsection{How it is verified}
\label{sec:security-privacy-verified}

\file{medorama/tests/privacy.rs} seeds a catalog with identifiers that cannot appear by chance (a
name, an MRN, patient, encounter and account numbers, a date of birth, staff names, an
institution, and real dates in the forms that tables, DICOM, MUSE, WFDB, EDF and EDF+ write them)
in every format the catalog serves: tables, notes, a WFDB record, GE MUSE XML, DICOM, an EDF+
file, a picture, a text file and an ICU recording. It requests every page's answers, every
entry's detail, every file and its parts, and every cohort and database export, as a user without
the PHI permission and as an administrator in Privacy Mode. It asserts that \textbf{no byte of any
answer holds any of them}, that no JSON number is a real time, that identifier searches are
refused or find nobody, that placeholders and pseudonyms appear, that unknown files are refused
and that pictures and DICOM files are marked. \file{tests/gateway.rs} checks the VitalStream
gateway with a real viewer, and \file{tests/connections.rs} repeats the check over a synthetic
OMOP database. Unit tests cover each de-identifier (\file{gate.rs}, \file{deid.rs},
\file{dicom.rs}, \file{pseudonyms.rs}, \file{gateway.rs}).

\section{Data at rest}
\label{sec:security-data}

\begin{itemize}
\item \textbf{The database} (SQLite, in WAL mode) lives in the catalog's home folder, outside the
  repository. The database and its \file{-wal} and \file{-shm} files are set to mode 0600, and
  the home folder to 0700 (\file{db.rs}).
\item \textbf{Encryption.} Builds with the \code{encryption} feature use SQLCipher (AES-256); the
  key comes from \code{\$MEDORAMA\_KEY} or from the file named by \code{key\_file}. Without it,
  rely on full-disk encryption.
\item \textbf{On PostgreSQL}, the database server encrypts at rest (an encrypted disk or volume;
  managed services such as Amazon RDS and Azure encrypt storage with their key services), and
  \code{key\_file} is refused. Medorama's account needs rights on its own database only.
\item \textbf{Identifiers} live in their own tables, so a dump of the clinical tables carries
  pseudonymous IDs only.
\item \textbf{Data sources are never written.} No part of Medorama writes into a data folder;
  indexes, results and converted files go to their own folders. \textbf{Connected databases are
  never written either}: Medorama's sessions there are read-only, whatever the account may do,
  and their data are not copied, apart from linked persons' identifiers.
\item \textbf{Index and results folders} hold derived data with real dates: protect them like the
  recordings.
\item \textbf{Secrets.} The catalog's secret (32 random bytes) is created once in the database's
  settings. The internal viewer's password file is mode 0600. Connection passwords entered in the
  page are written to \file{<home>/secrets/<name>.password} (mode 0600, in a folder of mode 0700),
  never to the database, and deleted when the database is disconnected. No secret is in the
  repository or the binary.
\item \textbf{Backups} of the database hold identifiers, the pseudonym mapping and the audit log:
  encrypt them, and restrict who may restore them.
\end{itemize}

\section{Data in transit}
\label{sec:security-transit}

\begin{itemize}
\item Medorama speaks plain HTTP/1.1 and listens on \textbf{127.0.0.1:8090} by default. Bound to
  any other address, or signing users in through a proxy, without \code{secure\_cookies = true},
  it prints a warning at start.
\item \textbf{TLS} is terminated by a reverse proxy in front of it (with
  \code{secure\_cookies = true}), or users reach it through an SSH tunnel
  (\cmd{ssh -N -L 8090:127.0.0.1:8090 user@host}). In proxy sign-in mode, the authenticating
  proxy terminates TLS.
\item The internal VitalStream viewer is reached over loopback only, with requests signed by its
  private key.
\item \textbf{PostgreSQL} (the catalog and connected databases) is reached over TLS through rustls
  with the system's certificate authorities (and \code{database\_ca\_file}, or a connection's
  \opt{-{}-ca-file}); the server's certificate and name are always verified. A server on another
  host must offer TLS unless its URL says \code{sslmode=disable}: libpq's default,
  \code{prefer}, would fall back to plain text, and Medorama requires TLS instead. A server on
  the same host may do without it (\file{sql/pg.rs}, \code{connect}).
\item Nothing is sent to any host other than those databases.
\end{itemize}

\section{Web application security}
\label{sec:security-web}

\begin{itemize}
\item \textbf{Content Security Policy} on every page (\code{guard}). For Medorama's page it is
  \code{default-src 'self'}, \code{frame-ancestors 'none'}, \code{base-uri 'none'},
  \code{form-action 'self'}. For the ECG and PACS viewers' pages it is \code{default-src 'self'},
  \code{blob:} only for images, media and frames, \code{object-src 'none'},
  \code{base-uri 'none'}, \code{form-action 'none'}, \code{frame-ancestors 'self'}. For the
  VitalStream viewer's page under \file{/viewer/} it is \code{default-src 'self'},
  \code{script-src 'self'}, \code{img-src 'self' data:}, \code{object-src 'none'},
  \code{base-uri 'none'}, \code{frame-ancestors 'self'}. No inline
  scripts run. The guard also sets \code{X-Frame-Options}, \code{X-Content-Type-Options: nosniff},
  \code{Referrer-Policy: no-referrer}, and \code{Cache-Control: no-store} on the API.
\item \textbf{CSRF.} Every API request that changes something must carry \code{X-Medorama: 1},
  which a cross-site form cannot set and a cross-site script cannot send without a CORS
  permission that Medorama never grants. Requests to the VitalStream viewer through the gateway
  that change something must carry \code{X-Medorama: 1} too (\file{server.rs} \code{guard}; the
  viewer's page sends it) and be JSON (\code{Content-Type: application/json}), which a cross-site
  page cannot send without the browser asking first. In password mode, session cookies are also
  \code{SameSite=Strict}.
\item \textbf{Active content in files.} SVG and XML files, and de-identified text files, are
  served with a \code{sandbox} Content Security Policy; the viewers show pictures, PDFs and videos
  from \code{blob:} addresses.
\item \textbf{Output escaping.} The pages build HTML only through an escaping helper
  (\code{esc}); there is no templating of untrusted HTML.
\item \textbf{No third-party code in the page.} The page and the viewers are plain JavaScript
  modules embedded in the binary; nothing loads from a CDN.
\end{itemize}

\section{Untrusted input}
\label{sec:security-input}

Every file and table is treated as untrusted.
\begin{itemize}
\item \textbf{Parsers} for tables, WFDB records, spreadsheets, GE MUSE XML, DICOM and EDF check
  every length against the data and never read out of bounds; a malformed file ends in an error,
  not a crash. The DICOM rewriter's tests truncate a file at every seventh byte and require an
  error or a result, never a panic.
\item \textbf{EDF and BDF} (\file{libs/engine/src/sources/edf.rs}, \code{header}): the header's
  size must equal 256 bytes per signal plus one; negative sample counts and non-positive record
  durations are refused; and the number of data records is taken from the file's size, never
  beyond it, whatever the header declares. Tests in \file{libs/engine/tests/edf.rs} include files
  that are not EDF. The Privacy Mode rewrite holds at most 64\,MB in memory.
\item \textbf{Fuzz-style tests} (\file{medorama/tests/robustness.rs}) feed the expression parser
  and the table importers broken and adversarial input.
\item \textbf{Coverage-guided fuzzing} (\file{medorama/fuzz/}, cargo-fuzz with libFuzzer) covers
  tables (CSV, spreadsheets), WFDB headers, GE MUSE XML, DICOM, EDF and BDF files (their header,
  their index, their identifiers and Privacy Mode's copy) and what users type (cohort
  expressions, searches, dates, numbers, units). CI builds every target on each change and runs
  each for 5 minutes every week (\cmd{make fuzz} locally).
\item \textbf{Paths.} \code{..} is refused in every served path; files are resolved after symbolic
  links and their real paths checked against the real paths of \code{data\_roots} (a link inside
  a root that leads out of it is outside). An attached folder's symbolic links are neither listed
  nor sent.
\item \textbf{DICOM sequences} nested more than 32 deep are refused (the rewrite would otherwise
  recurse without bound on a crafted file).
\item \textbf{Uploads} (imports) are limited to 1\,GiB and to administrators.
\item \textbf{The VitalStream viewers} read WFDB records with wfdb-x, a reimplementation of the
  WFDB C library that finds files only where they are named, keeps a header's signal files inside
  its folder, and has network access off. Medorama's fuzzing does not cover it: open WFDB records
  from trusted sources only. A folder of records is read within limits (5000 header files, 20
  folder levels, 10 seconds to count them), so a folder opened by mistake cannot exhaust the
  viewer.
\item \textbf{The viewers alone} never open anything outside their root, which is fixed when they
  are built; symbolic links are resolved, and a link that leads outside is neither listed nor
  opened. A build without password support refuses to run under Medorama.
\end{itemize}

\section{Audit logging}
\label{sec:security-audit}

The audit log (\file{db.rs} \code{audit}; \ui{Admin} $\rightarrow$ \ui{Audit}; \code{GET audit})
records, for each entry, the time, the user, the event, its target and a detail. Its events are
listed in \cref{tab:users-privacy-events}. In Privacy Mode, the audit page shows no target or
detail, because a search's text may name a patient. Medorama never deletes entries. An entry the
database cannot take at once (busy, or failing) is never lost: it is kept in order and written
before the next entry or by the next seal (\code{db::flush\_audit}), appended at once to
\code{audit\_export} when that is set (a line with \code{"pending": true}), and the server's
output says so.

\textbf{Tamper evidence} (\file{audit.rs}):
\begin{itemize}
\item \textbf{Seals.} Every minute, the server seals the entries written since the last seal. A
  seal is the SHA-256 chain of those entries, in order, starting from the previous seal's hash,
  and each entry records its seal. Entries committed late (in a long transaction) go into the
  next seal.
\item \textbf{Verification.} \cmd{medorama audit -{}-verify} (or \code{GET audit/verify}, for
  administrators) computes every chain again. An entry changed, deleted or added to a sealed
  batch, or a seal changed or deleted, shows as the first seal that no longer holds. Entries and
  seals are numbered by the database without reuse, so verification also counts what it numbered
  against what is left: the last seal deleted with its entries, or entries deleted before they
  were sealed, show as missing (on PostgreSQL, an entry lost with a transaction that failed shows
  the same way).
\item \textbf{Out of reach.} Whoever controls the database could rewrite entries and seals
  together, and its counters with them. Each seal and its entries are therefore also appended to
  \code{audit\_export} (a JSON-lines file, readable by its owner only, synced to disk). Put it on
  storage that the database's administrators cannot change (a write-once volume, a log
  collector). \code{GET audit/verify} compares the database's seals with it (each exported seal
  must be in the database, with the same hash): \textbf{the end of the log is protected by the
  export}, up to its last seal; entries written since are protected by the database's numbering
  only.
\end{itemize}

\section{Background jobs and plug-in actions}
\label{sec:security-jobs}

\begin{itemize}
\item \textbf{Jobs} (indexing, derived series, plug-ins) run in worker threads of the server, or
  in \cmd{medorama work} processes, from a queue in the database (\file{jobs.rs}).
\item \textbf{Plug-ins} are programs that an administrator installs (a folder with
  \file{action.toml}; \cref{chap:actions}). Each run is a child process in its own process group,
  with a timeout; cancelling or timing out kills the whole group. It receives a job description:
  the recording's path, its index folder, the encounters' and patients' \textbf{pseudonyms}
  (never their imported IDs), and an output folder.
\item \textbf{Not the server's secrets.} A plug-in does not inherit Medorama's own environment
  variables (\code{MEDORAMA\_\dots}: the database's URL and password, the encryption key), nor any
  whose name marks it as secret (containing \code{PASSWORD}, \code{SECRET}, \code{TOKEN},
  \code{PASSPHRASE} or \code{CREDENTIAL}, ending in \code{\_KEY}, PostgreSQL's
  \code{PGPASS\dots}, \code{DATABASE\_URL}). The rest stays, so that the plug-in finds its
  interpreter and libraries (\file{actions.rs}, \code{secret\_env}). It receives
  \code{MEDORAMA\_JOB}, its job's description.
\item \textbf{Trust.} A plug-in runs with the catalog's operating-system rights: it is trusted
  code, like the server. Install only plug-ins you have reviewed. Isolation is by process, not a
  sandbox: a plug-in can still read what the catalog's account can, including the database file.
\item \textbf{Adapters} (add-on readers of licensed formats; \cref{chap:addons}) are installed by
  an administrator. Their environment is cleared except for a short allow-list (\code{PATH},
  \code{HOME}, the temporary folders, the locale), every command has a time limit, and their
  standard error never reaches the program's log (\cref{sec:addons-limits}). They return
  identifiers only to Medorama, in answer to \code{identify}.
\end{itemize}

\section{Supply chain, build and repository}
\label{sec:security-supply}

\begin{itemize}
\item \textbf{Rust} is built with \cmd{cargo build -{}-locked}: dependency versions are pinned by
  \file{Cargo.lock} and can be vendored for offline builds (\cref{chap:installation}). Every
  third-party component and its license is listed in \cref{chap:licences}.
\item \textbf{The page and the viewers} are plain HTML, CSS and JavaScript in the repository,
  embedded in the binary: there is no package manager, no build step, and nothing fetched at run
  time.
\item \textbf{CI} runs formatting, lints (\cmd{clippy}) and every test on Linux, macOS and
  Windows, the same tests on PostgreSQL 14 and 17, a browser test of the page (headless Chrome),
  the encrypted build and the installers. Every workflow reads the repository only, its actions
  are pinned by commit, and every job has a time limit.
\item \textbf{Releases} are built by \file{.github/workflows/release.yml} on a version tag. The
  workflow runs the tests and \cmd{cargo deny} first, builds with a pinned Rust, packs
  reproducible archives with \file{THIRD\_PARTY\_CRATES.txt} (the licenses of every crate
  compiled in, written by \cmd{cargo about}) and wfdb-x's license (\cref{sec:licences-wfdbx}),
  and attaches them to a draft release with \file{SHA256SUMS} and a signed
  build provenance.
\item \textbf{Dependencies are checked} by \cmd{cargo deny} (\file{deny.toml}, for the whole Rust
  workspace; \cmd{make audit}) on every change and every week. Any RustSec advisory (vulnerable,
  unmaintained or yanked crates), any license not on the allowed list, or any crate from
  elsewhere than crates.io fails the build.
\item \textbf{The API is described} in \file{medorama/openapi.yaml} (OpenAPI 3.1): a test fails
  when a route is added, removed or changed on one side only.
\item \textbf{The repository holds no patient data.} The privacy check
  \file{tools/git-hooks/check-files}, run by the pre-commit and pre-merge-commit hooks on what
  is staged, refuses recordings, data files, DICOM and MUSE XML files (by their content),
  notebooks with outputs, real recording names and anything in the folders kept for local data.
  CI's \code{privacy} job (\file{.github/workflows/checks.yml}) runs the same check on every
  commit of each pull request and push, and on every file, so a commit made without the hooks is
  still caught before it is merged. Tests and demonstrations use synthetic data only, or public
  data sets kept outside git (\cref{sec:protecting-repository}).
\end{itemize}

\section{Regulatory mapping}
\label{sec:security-regulatory}

Medorama is a tool; compliance is a property of a deployment and its governance. This section
maps Medorama's behavior to common requirements, so that a privacy or security officer can judge
how it supports them. Medorama is designed to support these safeguards; it does not by itself
make a deployment compliant, and the deploying institution remains responsible.

\paragraph{HIPAA Security Rule.} \Cref{tab:security-hipaa-security} covers the technical and
selected administrative safeguards of 45 CFR Part 164, Subpart C.

\begin{xltabular}{\linewidth}{@{}P{4.2cm}P{5.2cm}L@{}}
\caption{HIPAA Security Rule safeguards, and how Medorama supports them.}
\label{tab:security-hipaa-security}\\
\toprule
\tabhead Safeguard & \tabhead How Medorama supports it & \tabhead The institution provides \\
\midrule
\endfirsthead
\toprule
\tabhead Safeguard & \tabhead How Medorama supports it & \tabhead The institution provides \\
\midrule
\endhead
\bottomrule
\endfoot
Access control, unique user identification (\S164.312(a)(2)(i)) & Personal accounts; in proxy
  mode, the institution's identities & Account provisioning and review \\
Access control, automatic logoff (\S164.312(a)(2)(iii)) & \code{session\_idle\_minutes},
  \code{session\_hours} (password mode) & The proxy's or identity provider's session limits
  (proxy mode) \\
Access control, encryption and decryption (\S164.312(a)(2)(iv)) & Optional SQLCipher encryption
  of the SQLite catalog & Full-disk and backup encryption; PostgreSQL storage encryption \\
Audit controls (\S164.312(b)) & The audit log, sealed every minute, exportable
  (\cref{sec:security-audit}) & Regular review; storage for \code{audit\_export} \\
Integrity (\S164.312(c)(1)) & Data sources and connected databases are only read; audit seals
  detect changes to the log & Backups and their integrity \\
Person or entity authentication (\S164.312(d)) & Argon2id passwords with throttling; proxy
  sign-in for SSO and MFA & The identity provider and MFA policy \\
Transmission security (\S164.312(e)(1)) & Loopback by default; verified TLS to PostgreSQL & TLS
  termination at a reverse proxy, network policy \\
Information access management (\S164.308(a)(4)) & Roles, the PHI permission, Privacy Mode by
  default & Deciding who receives the PHI permission \\
Log-in monitoring (\S164.308(a)(5)(ii)(C)) & \code{login} and \code{login\_failed} entries &
  Monitoring them \\
Risk analysis, contingency plan, physical safeguards (\S164.308, \S164.310) & This chapter, as
  input to the analysis & The analysis, backups, disaster recovery, facility controls \\
\end{xltabular}

\paragraph{HIPAA Privacy Rule, de-identification.} \Cref{tab:security-hipaa} lists the 18 kinds
of identifiers of the Safe Harbor method (45 CFR \S164.514(b)(2)), and what Privacy Mode does with
each.

\begin{xltabular}{\linewidth}{@{}P{5.4cm}L@{}}
\caption{HIPAA's Safe Harbor identifiers, and Privacy Mode.}
\label{tab:security-hipaa}\\
\toprule
\tabhead Identifier & \tabhead Privacy Mode \\
\midrule
\endfirsthead
\toprule
\tabhead Identifier & \tabhead Privacy Mode \\
\midrule
\endhead
\bottomrule
\endfoot
Names & \code{[Patient Name]}; staff \code{[Staff Name]} \\
Geographic subdivisions smaller than a state & Not a catalog field: a column holding addresses
  is unrecognized, and Privacy Mode withholds it (\code{[Withheld]}) \\
All elements of dates (except year) related to an individual & \textbf{Moved back, not removed}
  (see below) \\
Telephone and fax numbers, e-mail addresses & Likewise unrecognized columns, withheld \\
Social security numbers & Not held by design; an unrecognized column holding them is withheld
  (only measurement columns of unrecognized data are shown) \\
Medical record numbers & \code{[MRN]} \\
Health plan beneficiary numbers, account numbers & \code{[CSN]} \\
Certificate or license numbers, vehicle and device identifiers & DICOM device serial numbers
  removed; other such columns unrecognized, so withheld \\
Web addresses, IP addresses & Not held \\
Biometric identifiers, full-face photographs & \textbf{Not removed}: pictures and images are sent
  as they are, marked \\
Any other unique identifying number or code & Imported patient and encounter IDs replaced by
  pseudonyms; DICOM UIDs replaced \\
Ages over 89 & \code{90+} \\
\end{xltabular}

\paragraph{Dates are shifted, not removed.} Safe Harbor removes every element of a date except
the year. Privacy Mode keeps dates, moved by an offset per patient, because intervals are the
substance of research. Data viewed or exported in Privacy Mode are therefore \textbf{not}
de-identified under Safe Harbor. They are suited to use as a \textbf{limited data set} with a data
use agreement (\S164.514(e)), or under an \textbf{expert determination} (\S164.514(b)(1)) that
weighs the offset, the residual quasi-identifiers and the recipients. Images and free text may
still identify.

\paragraph{GDPR.} Data shown or exported in Privacy Mode are pseudonymized data (Art.~4(5)) and
remain personal data. Medorama is designed to support data protection by design and by default
(Art.~25: Privacy Mode is the default for every user but administrators) and the security of
processing (Art.~32: pseudonymization, access control, optional encryption, audit logging). The
pseudonym mapping and the identifiers are kept separately, in the catalog's own tables, under the
catalog's access controls. The controller's other duties (lawful basis, records of processing,
data subjects' rights, retention) remain with the institution.

\section{Deployment checklist}
\label{sec:security-checklist}

Go through this list before real data come in. How to carry out each step is described in
\cref{chap:deployment} and \cref{chap:sign-in}.
\begin{enumerate}
\item A dedicated host or virtual machine, an operating-system account for the catalog, full-disk
  encryption, and the operating system kept up to date. No untrusted local accounts on the host.
\item The catalog's home folder on a local disk (SQLite's locking is unreliable on network file
  systems), owned by the catalog's account.
\item Build with \code{encryption} and keep the key outside the home folder
  (\code{\$MEDORAMA\_KEY} from the service manager, or a \code{key\_file} readable only by the
  account), or rely on disk encryption.
\item Keep \code{host = "127.0.0.1"}. Put a reverse proxy with TLS in front, and set
  \code{secure\_cookies = true}, or give users SSH tunnels. Never expose port 8090 directly.
\item Where policy requires single sign-on or MFA, use proxy sign-in (\cref{chap:sign-in}):
  \begin{itemize}
  \item set \code{sign\_in = "proxy"}, \code{proxy\_user\_header}, and \code{proxy\_addresses}
    to the proxy's address only;
  \item make Medorama reachable only from the proxy (loopback, or a firewall rule);
  \item make the proxy remove any incoming copy of the header, and verify it with a forged
    request (step 6 of \cref{sec:sign-in-setup});
  \item create accounts under the names the proxy sends (\cmd{medorama user add NAME -{}-role R
    -{}-no-password}), and set \code{proxy\_logout\_url};
  \item set session lifetime and idle timeout at the proxy or the identity provider.
  \end{itemize}
\item Set \code{data\_roots} to the folders that hold the data, and nothing more; mount them
  read-only for the catalog's account where possible.
\item Grant the PHI permission only to the people who need identifiers; everyone else works in
  Privacy Mode. Review the list of users regularly, and disable accounts promptly when people
  leave.
\item In \ui{Admin} $\rightarrow$ \ui{Privacy Mode}, review what is withheld and choose the date
  offset and each patient's days; record them under your governance, and do not disclose the
  offset.
\item Install only reviewed plug-ins and add-ons.
\item Back up the database (it holds identifiers, the pseudonym mapping and the audit log),
  encrypt the backups, and test restoring them.
\item Set \code{audit\_export} to a file on write-once storage, or one that your log management
  collects; run \cmd{medorama audit -{}-verify} regularly (it fails when the log was altered).
\item Run Medorama as a service: a systemd user service, a launchd agent, or on Windows a script
  in the Startup folder (\cref{sec:deployment-service}).
\item \textbf{On PostgreSQL:} a database and an account for Medorama alone; TLS on the server
  (Medorama requires it, and verifies the certificate and the name); the password in a 0600 file
  or the service's environment; the server's storage encrypted; \cmd{pg\_dump} backups encrypted.
\item \textbf{Connected databases:} an account that can only read, on the views' tables only; the
  network path allowed from Medorama's hosts only; TLS; a time limit and a row limit suited to the
  database (\opt{-{}-timeout-s}, \opt{-{}-max-rows}).
\item \textbf{No demo with real patients:} install without \opt{-{}-demo} (the default), so that
  \cmd{init} creates an empty catalog. If a catalog holds the demo (a cohort named DEMO), an
  administrator deletes that cohort before real data arrive; its invented patients then leave the
  catalog.
\end{enumerate}

\section{Known limitations and residual risks}
\label{sec:security-limitations}

The work planned against these limitations is listed in the repository's \file{docs/todo.md}.
\begin{itemize}
\item \textbf{Pixels, documents, video, free text.} Names burned into images, scanned PDFs and
  video frames, and free text that names relatives or staff, are not removed; they are marked.
  Free-text de-identification and optical character recognition of burned-in text are out of
  scope for now.
\item \textbf{EDF+ annotations in VitalStream} are free text. In Privacy Mode they are treated
  like notes: withheld when the administrator withholds notes, otherwise shown with the patient's
  own names and date of birth replaced and the identifiers the catalog knows replaced. Other
  names in them are not removed. In files sent to the ECG viewer, annotation texts are always
  removed.
\item \textbf{Quasi-identifiers.} Sex, a capped age, rare diagnoses, units and beds, small cohorts
  and exact intervals can single out a person when combined with outside knowledge.
\item \textbf{Date offsets.} Learning one real date of a patient reveals that patient's offset,
  and every other patient's to within \code{spread\_days}.
\item \textbf{Every user sees every patient} of a catalog, as their mode allows: there are no
  per-patient or per-study access lists.
\item \textbf{Proxy sign-in trusts the proxy and the host.} Anyone who can connect to Medorama's
  port from an address in \code{proxy\_addresses}, including any process on the host with the
  default addresses, can name any user. Roles and the PHI permission are not taken from the
  identity provider.
\item \textbf{Password sign-in has no multi-factor authentication} and no password expiry. Use
  proxy sign-in where policy requires them.
\item \textbf{Administrators are trusted.} The audit log lives in the same database. Its seals
  show a change after the fact, and its export (\code{audit\_export}) keeps a copy out of the
  database's reach, but a host administrator who changes both the database and the export file at
  once is not detected.
\item \textbf{Plug-ins and add-ons} run with the catalog's rights, without its secrets in their
  environment.
\item \textbf{Connected databases' own access rules} are not Medorama's: every Medorama user sees
  the data of linked persons, as their mode allows. Connect only databases that all Medorama
  users may use.
\item \textbf{A connected database that does not answer} leaves patient pages without its data
  (the banner names it) and makes cohort searches fail rather than silently miss it.
\item \textbf{Unsupported encodings are refused, not converted.} Deflated or big-endian DICOM, XML
  that is not GE MUSE, text attachments, and EDF+ files with annotations larger than 64\,MB are
  refused in Privacy Mode.
\end{itemize}

\section{Verifying these claims}
\label{sec:security-verifying}

From a clone of the repository, the following commands run the tests that check the controls
above. In \file{medorama/}:
\begin{lstlisting}[style=shell]
cargo test --locked                # every test, including those below
cargo test --test privacy          # Privacy Mode end to end: no identifier in any answer
cargo test --test privacy_search   # searches, filters and sorts in Privacy Mode: shown values only
cargo test --test signin           # password and proxy sign-in
cargo test --test gateway          # the VitalStream gateway with a real viewer
cargo test --test connections      # a connected OMOP database, Privacy Mode included
cargo test --test browser          # the page in headless Chrome, Privacy Mode on screen
cargo test --test openapi          # the API's description names every route
cargo test --lib -- gate:: deid:: dicom:: pseudonyms:: config::   # each de-identifier
cargo clippy --all-targets         # lints
MEDORAMA_TEST_POSTGRES=postgres://user@host/db cargo test   # again, the catalog in PostgreSQL
medorama audit --verify            # the audit log's seals hold (on a deployed catalog)
\end{lstlisting}
At the repository's root:
\begin{lstlisting}[style=shell]
cargo test -p recording-engine --test edf    # the EDF and BDF reader
cargo deny check                             # dependencies: advisories, licenses, sources
make fuzz                                    # the parsers, fuzzed (EDF and BDF included)
tools/git-hooks/check-files --tree HEAD      # no patient data in any file of the repository
\end{lstlisting}

\begin{xltabular}{\linewidth}{@{}P{5.6cm}L@{}}
\caption{Where to inspect each control in the code.}
\label{tab:security-inspect}\\
\toprule
\tabhead To inspect & \tabhead Read \\
\midrule
\endfirsthead
\toprule
\tabhead To inspect & \tabhead Read \\
\midrule
\endhead
\bottomrule
\endfoot
Sign-in: methods, the proxy's header, password sign-in and throttling &
  \file{medorama/src/signin.rs}; the settings in \file{medorama/src/config.rs}
  (\code{sign\_in\_method}) \\
Sessions, the Privacy Mode decision, the last layer, security headers &
  \file{medorama/src/server.rs}: \code{session\_user}, \code{in\_privacy}, \code{privacy},
  \code{guard} \\
The gate: dictionary, placeholders, files, table cells & \file{medorama/src/gate.rs} \\
Dates, ages, WFDB headers, MUSE XML, EDF headers and annotations & \file{medorama/src/deid.rs} \\
EDF+ annotations in VitalStream & \file{medorama/src/gateway.rs}: \code{gate\_annotations} \\
DICOM & \file{medorama/src/dicom.rs} \\
Files in Privacy Mode & \file{medorama/src/server.rs}: \code{attachment\_response},
  \code{edf\_deidentified}, \code{file\_response} \\
Pseudonyms & \file{medorama/src/pseudonyms.rs}; the \code{pseudonym} column of \code{patients}
  and \code{encounters}, \code{pseudonym\_aliases} (\file{medorama/src/schema/}) \\
Accounts and sessions & \file{medorama/src/users.rs} \\
The VitalStream gateway and the viewer's supervisor & \file{medorama/src/gateway.rs},
  \file{medorama/src/viewer.rs} \\
The database, migrations, file permissions, encryption key & \file{medorama/src/db.rs},
  \file{medorama/src/schema/} \\
SQLite and PostgreSQL behind one interface, TLS & \file{medorama/src/sql/} \\
Connected databases: read-only access, views, the registry & \file{medorama/src/sources/} \\
The audit log's seals & \file{medorama/src/audit.rs}; the tables \code{audit} and
  \code{audit\_seals} \\
Plug-ins' environment & \file{medorama/src/actions.rs} (\code{secret\_env}),
  \file{medorama/src/jobs.rs} \\
The API, endpoint by endpoint & \file{medorama/openapi.yaml} \\
Dependency policy & \file{deny.toml} (the Rust workspace: Medorama, the libraries,
  \code{vitalstream-rs}) \\
\end{xltabular}
